% =====================================================================
% 4. DEVELOPING MACHINE LEARNING MODELS
% =====================================================================
\section{Developing Machine Learning models}

\subsection{Model Selection}

Each model is defined in \texttt{src/models.py} as a \texttt{Pipeline} with two steps,
\texttt{prep} (a \texttt{Column\allowbreak{}Transformer} from \texttt{src/\allowbreak{}pipelines.\allowbreak{}py}) and \texttt{model}.
Linear models and support vector machines receive standardised inputs and one-hot encoded race
(\texttt{handle\_\allowbreak{}unknown="ignore"}, needed because one race category has a single patient and may be
absent from a training fold); trees and ensembles receive unscaled inputs. Naive Bayes has its own
preprocessing: age is discretised into quantile bins (\texttt{KBinsDiscretizer}) so that a
\texttt{BernoulliNB} can model every column. Tables~\ref{tab:mllb-reg-grids} and~\ref{tab:mllb-clf-grids}
list the models and the grids searched; the grids are small (at most 27 configurations) because
nested cross-validation multiplies the number of fits by 30.

\begin{table}[H]
\centering
\caption{Regression models and hyperparameter grids.}
\label{tab:mllb-reg-grids}
\resizebox{\textwidth}{!}{\input{tables/reg_models_grids.tex}}
\end{table}

\begin{table}[H]
\centering
\caption{Classification models and hyperparameter grids (the SVM grid is a list of two sub-grids,
linear and RBF).}
\label{tab:mllb-clf-grids}
\resizebox{\textwidth}{!}{\input{tables/clf_models_grids.tex}}
\end{table}

\subsection{Model Tuning}

Tuning is done by \texttt{GridSearchCV} with 3 inner folds (stratified for classification,
\texttt{n\_jobs=-1}), scoring \texttt{roc\_auc} for classification and
\texttt{neg\_\allowbreak{}root\_\allowbreak{}mean\_\allowbreak{}squared\_\allowbreak{}error} for regression, inside each of the 10 outer folds
(\texttt{evaluation.\allowbreak{}nested\_\allowbreak{}cv}). Because each outer fold runs its own search, the chosen values can
be compared across folds: a value chosen in most folds indicates a stable choice
(Tables~\ref{tab:mllb-reg-hp} and~\ref{tab:mllb-clf-hp}). The final models used for interpretation are
obtained by \texttt{evaluation.\allowbreak{}fit\_\allowbreak{}final}, which runs the same search once on all patients; its result
is the last column. For logistic regression the final search chose $C=\LRFinalC$ although the folds
chose smaller values; Table~\ref{tab:mllb-flat} shows why: \LRFlatWithin{} of the \LRFlatTotal{}
configurations are within 0.005 AUC of the best one.

\begin{table}[H]
\centering
\caption{Regression: hyperparameters chosen in the 10 outer folds (count in brackets) and in the final
model.}
\label{tab:mllb-reg-hp}
\resizebox{\textwidth}{!}{\input{tables/reg_hyperparameters.tex}}
\end{table}

\begin{table}[H]
\centering
\caption{Classification: hyperparameters chosen in the 10 outer folds and in the final model.}
\label{tab:mllb-clf-hp}
\resizebox{\textwidth}{!}{\input{tables/clf_hyperparameters.tex}}
\end{table}

\begin{table}[H]
\centering
\caption{Flatness of the inner-CV surface of the final classification searches.}
\label{tab:mllb-flat}
\resizebox{\textwidth}{!}{\input{tables/clf_search_flatness.tex}}
\end{table}

A reproducibility problem was found at this stage: two runs of the classification notebook chose
different hyperparameters for logistic regression, because the \texttt{liblinear} solver shuffles the
data and its \texttt{random\_state} had not been fixed. After fixing it, two complete runs give
identical out-of-fold probabilities, and the author's run reproduces every value of the JSON summaries
(Appendix A).

\subsection{Model Evaluation}

\textbf{Regression.} Figure~\ref{fig:mllb-reg-comp} shows the out-of-fold MAE and $R^2$ of every model
with bootstrap intervals (all models between \RegMinMAE{} and \RegMaxMAE{} years of MAE, against
\RegBaseMAE{} for the mean baseline), and Table~\ref{tab:mllb-reg-paired} the paired differences. The
multiple linear regression was refitted with statsmodels for its diagnostics
(Figure~\ref{fig:mllb-ols}); one patient with leverage 1 (the only patient in one race category) is
excluded from the plots because its residual is zero by construction. Figure~\ref{fig:mllb-lasso} shows
the Lasso coefficient path with the chosen penalty ($\alpha=\LassoAlpha$, \LassoNZero{} of
\LassoNCoef{} coefficients set to zero), Figure~\ref{fig:mllb-regtree} the final regression tree
(depth \RegTreeDepth, \RegTreeLeaves{} leaves) and Figure~\ref{fig:mllb-regimp} the permutation
importance computed on the outer test folds.

\begin{figure}[H]
\centering
\includegraphics[width=\textwidth]{reg_model_comparison.png}
\caption{Regression: out-of-fold MAE and $R^2$ with 95\% bootstrap intervals.}
\label{fig:mllb-reg-comp}
\end{figure}

\begin{table}[H]
\centering
\caption{Regression: paired bootstrap of the MAE difference (years) against the mean baseline and
against the model with the lowest MAE.}
\label{tab:mllb-reg-paired}
\small
\input{tables/reg_paired_differences.tex}
\end{table}

\begin{figure}[H]
\centering
\includegraphics[width=0.85\textwidth]{reg_ols_diagnostics.png}
\caption{Residual diagnostics of the multiple linear regression (all patients).}
\label{fig:mllb-ols}
\end{figure}

\begin{figure}[H]
\centering
\begin{minipage}[t]{0.49\textwidth}\centering
\includegraphics[width=\textwidth]{reg_lasso_path.png}
\caption{Lasso path (dotted line: chosen $\alpha$).}
\label{fig:mllb-lasso}
\end{minipage}\hfill
\begin{minipage}[t]{0.49\textwidth}\centering
\includegraphics[width=\textwidth]{reg_permutation_importance.png}
\caption{Regression: permutation importance on the outer test folds.}
\label{fig:mllb-regimp}
\end{minipage}
\end{figure}

\begin{figure}[H]
\centering
\includegraphics[width=\textwidth]{reg_decision_tree.png}
\caption{Final regression tree (node colour: darker = older predicted age).}
\label{fig:mllb-regtree}
\end{figure}

\textbf{Classification.} Figure~\ref{fig:mllb-clf-comp} shows the out-of-fold AUC and balanced
accuracy, Table~\ref{tab:mllb-clf-acc} the accuracy and the per-fold AUC, and
Table~\ref{tab:mllb-clf-paired} the paired AUC differences with the best model. The confusion matrices
and the sensitivity--specificity trade-off as the threshold moves are in Figure~\ref{fig:mllb-confusion};
the threshold is left at 0.5, because tuning it on the same predictions would bias the reported
metrics. Figure~\ref{fig:mllb-clftree} shows the final pruned tree (\texttt{ccp\_alpha} $=\ClfTreeAlpha$,
\ClfTreeCriterion{} criterion), recoloured with the project colours (sklearn's default paints LGG in
orange), and Figures~\ref{fig:mllb-logit}--\ref{fig:mllb-clfimp} the logistic coefficients and the
permutation importance.

\begin{figure}[H]
\centering
\includegraphics[width=\textwidth]{clf_model_comparison.png}
\caption{Classification: out-of-fold AUC and balanced accuracy with 95\% bootstrap intervals.}
\label{fig:mllb-clf-comp}
\end{figure}

\begin{table}[H]
\centering
\begin{minipage}[t]{0.52\textwidth}\centering
\caption{Accuracy (95\% CI) and AUC per outer fold (mean $\pm$ SD).}
\label{tab:mllb-clf-acc}
\resizebox{\textwidth}{!}{\input{tables/clf_accuracy_folds.tex}}
\end{minipage}\hfill
\begin{minipage}[t]{0.45\textwidth}\centering
\caption{Paired bootstrap of the AUC difference with the best model.}
\label{tab:mllb-clf-paired}
\resizebox{\textwidth}{!}{\input{tables/clf_paired_differences.tex}}
\end{minipage}
\end{table}

\begin{figure}[H]
\centering
\includegraphics[width=\textwidth]{clf_confusion_threshold.png}
\caption{Out-of-fold confusion matrices at threshold 0.5 and sensitivity/specificity as a function of
the threshold for the model with the highest AUC.}
\label{fig:mllb-confusion}
\end{figure}

\begin{figure}[H]
\centering
\includegraphics[width=0.9\textwidth]{clf_decision_tree.png}
\caption{Final pruned classification tree (blue = LGG, orange = GBM; darker = purer node).}
\label{fig:mllb-clftree}
\end{figure}

\begin{figure}[H]
\centering
\begin{minipage}[t]{0.49\textwidth}\centering
\includegraphics[width=\textwidth]{clf_logistic_coefficients.png}
\caption{Final logistic regression: coefficients on standardised inputs.}
\label{fig:mllb-logit}
\end{minipage}\hfill
\begin{minipage}[t]{0.49\textwidth}\centering
\includegraphics[width=\textwidth]{clf_permutation_importance.png}
\caption{Classification: permutation importance on the outer test folds.}
\label{fig:mllb-clfimp}
\end{minipage}
\end{figure}

\textbf{Multiclass classification (notebook 06).} The target is built by \texttt{data.\allowbreak{}diagnosis\_\allowbreak{}family},
which maps the six TCGA diagnoses to four families with the dictionary \texttt{config.\allowbreak{}DIAGNOSIS\_\allowbreak{}FAMILY}
and raises an error if a diagnosis is not mapped (Table~\ref{tab:mllb-mc-classes}). The same
\texttt{nested\_cv} function is used with \texttt{task="multiclass"}: the folds are stratified by family,
and the out-of-fold prediction of every patient is a row of four probabilities. The inner criterion is
sklearn's \texttt{roc\_auc\_ovo} (Hand--Till macro AUC); for the 2000 bootstrap resamples the same
metric is recomputed by \texttt{evaluation.\allowbreak{}multiclass\_\allowbreak{}auc}, written with the rank-based binary AUC of the
project and checked to give exactly the value of \texttt{roc\_\allowbreak{}auc\_\allowbreak{}score(multi\_\allowbreak{}class="ovo")}. The
multinomial logistic regression uses the \texttt{saga} solver, the only one in scikit-learn~1.6 that
supports both penalties for the multinomial model. Table~\ref{tab:mllb-mc-grids} lists the models,
Table~\ref{tab:mllb-mc-hp} the chosen hyperparameters and Table~\ref{tab:mllb-mc-paired} the paired
comparison; Figures~\ref{fig:mllb-mc-comp} and~\ref{fig:mllb-mc-imp} show the metrics and the
permutation importance (drop in macro AUC).

\begin{table}[H]
\centering
\caption{Multiclass target: TCGA diagnoses grouped into histological families.}
\label{tab:mllb-mc-classes}
\small
\input{tables/mc_classes.tex}
\end{table}

\begin{table}[H]
\centering
\caption{Multiclass models and hyperparameter grids.}
\label{tab:mllb-mc-grids}
\resizebox{\textwidth}{!}{\input{tables/mc_models_grids.tex}}
\end{table}

\begin{table}[H]
\centering
\caption{Multiclass: hyperparameters chosen in the 10 outer folds and in the final model.}
\label{tab:mllb-mc-hp}
\resizebox{\textwidth}{!}{\input{tables/mc_hyperparameters.tex}}
\end{table}

\begin{figure}[H]
\centering
\includegraphics[width=\textwidth]{mc_model_comparison.png}
\caption{Multiclass: out-of-fold macro AUC and balanced accuracy with 95\% bootstrap intervals.}
\label{fig:mllb-mc-comp}
\end{figure}

\begin{table}[H]
\centering
\caption{Multiclass: paired bootstrap of the macro-AUC difference with the best model.}
\label{tab:mllb-mc-paired}
\small
\input{tables/mc_paired_differences.tex}
\end{table}

\begin{table}[H]
\centering
\caption{Per-family performance of the best multiclass model (out-of-fold; one-vs-rest AUC with 95\%
bootstrap interval).}
\label{tab:mllb-mc-perclass}
\footnotesize
\input{tables/mc_per_class.tex}
\end{table}

\begin{figure}[H]
\centering
\includegraphics[width=0.75\textwidth]{mc_permutation_importance.png}
\caption{Multiclass: permutation importance on the outer test folds.}
\label{fig:mllb-mc-imp}
\end{figure}

\textbf{SHAP values (notebook 07).} The notebook reads the JSON summaries of notebooks 05 and 06 to find
the winning models (the best one and those whose paired difference with it includes zero), refits them
with \texttt{fit\_final} and asserts that the hyperparameters are the ones saved by those notebooks.
\texttt{explain.\allowbreak{}shap\_\allowbreak{}values} applies the pipeline's preprocessing and chooses the explainer by model
type (\texttt{TreeExplainer} for forests, trees and XGBoost; \texttt{LinearExplainer} with all patients
as background for logistic regression); it returns the values, the base value, the units and the raw
model output, and \texttt{explain.\allowbreak{}additivity\_\allowbreak{}error} checks that base value plus contributions reproduces
every prediction. The summary plots (\texttt{explain.\allowbreak{}beeswarm}) are drawn with Matplotlib rather than
with the library's own plots, to keep the project style. Table~\ref{tab:mllb-shap-agree} compares the
SHAP ranking with the permutation importance, and Figure~\ref{fig:mllb-shap-age} shows the dependence of
the age contribution.

\begin{table}[H]
\centering
\caption{Agreement between SHAP (share of mean $|$SHAP$|$, race columns summed) and permutation importance.}
\label{tab:mllb-shap-agree}
\resizebox{\textwidth}{!}{\input{tables/shap_binary_agreement.tex}}
\end{table}

\begin{figure}[H]
\centering
\includegraphics[width=\textwidth]{shap_binary_age.png}
\caption{SHAP dependence of age for the three tied binary models (colour: IDH1 status).}
\label{fig:mllb-shap-age}
\end{figure}

\textbf{Causal analysis (notebook 08).} \texttt{causal.\allowbreak{}mediation\_\allowbreak{}decomposition} fits the three least-squares
regressions with \texttt{numpy.\allowbreak{}linalg.\allowbreak{}lstsq} (fast enough for 2000 bootstrap resamples) and returns the
identity check $c-(c'+ab)$, which the notebook asserts to be below $10^{-8}$;
\texttt{causal.\allowbreak{}ols\_\allowbreak{}table} refits the models with statsmodels and HC3 standard errors for the reported
intervals, because the residuals of age are heteroscedastic (Figure~\ref{fig:mllb-ols}). Race enters as a single
white/non-white indicator: with the four categories, the patient alone in one category has leverage 1,
and the HC3 correction divides by $1-h_{ii}$. Table~\ref{tab:mllb-causal-adj} lists the four adjustment
sets and Table~\ref{tab:mllb-causal-int} the interaction model.

\begin{table}[H]
\centering
\caption{Coefficient of IDH mutation on age at diagnosis (years) for different adjustment sets (HC3
intervals).}
\label{tab:mllb-causal-adj}
\small
\input{tables/causal_adjustment_sets.tex}
\end{table}

\begin{table}[H]
\centering
\caption{Age at diagnosis with an IDH $\times$ grade interaction (OLS, HC3 intervals).}
\label{tab:mllb-causal-int}
\small
\input{tables/causal_interaction.tex}
\end{table}
